\chapter{METHODOLOGY}%
\label{chap:methodology}

This chapter describes the methodology adopted for the design and implementation of the 
digital clock system. It explains the overall development approach, the system architecture, 
and the sequence of steps followed to achieve the project objectives.

\section{OVERVIEW OF PROPOSED METHODOLOGY}

The proposed methodology follows a systematic approach to develop a microcontroller-based 
digital clock using the AVR ATmega32. The design process begins with requirement analysis, 
followed by hardware and software design, implementation, and testing.

The system is developed using an interrupt-driven approach to ensure accurate time measurement. 
Hardware timers are configured to generate precise one-second intervals, and multiplexing is 
employed to control multiple seven-segment displays efficiently. Simulation using Proteus is 
performed to validate the design before final testing.

\section{WATERFALL DEVELOPMENT MODEL}

The Waterfall development model is adopted for this project due to its structured and 
sequential nature, which is suitable for small-scale embedded system applications with 
well-defined requirements.

The stages of the Waterfall model used in this project are as follows:
\begin{itemize}
    \item \textbf{Requirement Analysis}: Identification of functional and non-functional requirements of the digital clock system.
    \item \textbf{System Design}: Design of hardware connections and software logic, including timer configuration and display multiplexing.
    \item \textbf{Implementation}: Development of embedded C code using Microchip Studio and creation of the circuit in Proteus.
    \item \textbf{Testing}: Verification of correct time counting and display operation through simulation.
    \item \textbf{Maintenance}: Provision for future modifications and enhancements.
\end{itemize}
\enumitemEndSpacing{}

The Waterfall model ensures clear documentation at each stage and simplifies project management, 
making it appropriate for academic projects.

\section{ARCHITECTURE OF PROPOSED SYSTEM}

The architecture of the proposed digital clock system consists of three main components: 
the microcontroller unit, the display unit, and the power supply.

The AVR ATmega32 microcontroller serves as the core processing unit of the system. It is 
responsible for generating time intervals using Timer1 configured in Clear Timer on Compare (CTC) 
mode and handling interrupts to update time variables. The microcontroller also controls the 
multiplexing of four seven-segment displays through its I/O ports.

The display unit comprises four common-cathode seven-segment displays arranged to represent 
minutes and seconds in MM:SS format. Segment lines are connected to a common port through 
current-limiting resistors, while individual digit enable lines are used to activate each 
display sequentially.

The power supply provides a regulated 5V DC voltage required for the operation of the 
microcontroller and display unit. Figure~\ref{fig:system_architecture} illustrates the 
overall architecture of the proposed system.

\begin{figure}
\centering
\includegraphics[width=\figWidth{}]{system_architecture.png}
\caption{System Architecture of Digital Clock System}%
\label{fig:system_architecture}
\end{figure}
